\documentclass[a4paper,12pt]{article}

\usepackage[fontset=fandol]{ctex}
\usepackage{amsmath,amssymb}
\usepackage{graphicx}
\usepackage[margin=2.35cm]{geometry}
\usepackage[most]{tcolorbox}
\usepackage{etoolbox}
\usepackage{listings}
\usepackage{hyperref}
\usepackage{booktabs}
\usepackage{subcaption}
\usepackage{float}
\usepackage{xcolor}
\usepackage{caption}
\usepackage{tikz}
\IfFileExists{pgfplots.sty}{
  \usepackage{pgfplots}
  \pgfplotsset{compat=1.18}
}{}

\hypersetup{colorlinks=true,linkcolor=blue!60!black,urlcolor=blue!60!black}
\setlength{\parskip}{0.45em}
\setlength{\parindent}{0pt}
\captionsetup{font=small,labelfont=bf}

\newtcolorbox{knowledgebox}[1]{
  enhanced,
  colback=blue!5!white,
  colframe=blue!70!black,
  colbacktitle=blue!70!black,
  coltitle=white,
  fonttitle=\bfseries,
  title=#1,
  sharp corners
}

\newtcolorbox{importantbox}[1]{
  enhanced,
  colback=yellow!10!white,
  colframe=yellow!80!black,
  colbacktitle=yellow!80!black,
  coltitle=black,
  fonttitle=\bfseries,
  title=#1,
  sharp corners
}

\newtcolorbox{warningbox}[1]{
  enhanced,
  colback=red!5!white,
  colframe=red!70!black,
  colbacktitle=red!70!black,
  coltitle=white,
  fonttitle=\bfseries,
  title=#1,
  sharp corners
}

\lstset{
  basicstyle=\ttfamily\small,
  breaklines=true,
  frame=single,
  numbers=left,
  numberstyle=\tiny\color{gray},
  captionpos=b,
  extendedchars=false
}

% Fill these metadata fields during generation.
\newcommand{\notetitle}{课程讲义}
\newcommand{\notesubtitle}{}
\newcommand{\noteauthors}{Codex}
\newcommand{\notedate}{\today}
\newcommand{\videochannel}{}
\newcommand{\videopublishdate}{}
\newcommand{\videoduration}{}
\newcommand{\videourl}{}
\newcommand{\videocoverpath}{}
\newcommand{\inputmode}{}
\newcommand{\processingnote}{}

\newcommand{\sourceframefig}[4]{%
\begin{figure}[H]
\centering
\includegraphics[width=#4\textwidth]{#1}
\caption{#2\protect\footnotemark}
\end{figure}
\footnotetext{视频画面时间区间：#3。若为裁剪图，时间区间仍对应原始视频画面。}
}

\newcommand{\teachingdiagram}[4]{%
\begin{figure}[H]
\centering
\includegraphics[width=#4\textwidth]{#1}
\caption{#2。教学图来源依据：#3。此图为重写生成的示意图，不是原始视频截图。}
\end{figure}
}

\begin{document}

\begin{titlepage}
\centering
{\Large 课程讲义\par}
\vspace{1.0cm}
{\huge\bfseries \notetitle\par}
\ifdefempty{\notesubtitle}{}{\vspace{0.35cm}{\Large \notesubtitle\par}}
\vspace{0.65cm}
{\large \noteauthors\par}
\vspace{0.25cm}
{\large \notedate\par}
\vspace{0.9cm}

\ifdefempty{\videocoverpath}{
  {\small 未提供可用封面图。生成最终版前应补充平台封面或本地预览帧。\par}
}{
  \includegraphics[width=0.86\textwidth,height=0.43\textheight,keepaspectratio]{\videocoverpath}\par
}

\vfill
\begin{tcolorbox}[width=0.9\textwidth,colback=black!2!white,colframe=black!60,sharp corners]
\textbf{视频作者/频道：}\videochannel\par
\textbf{发布日期：}\videopublishdate\par
\textbf{视频时长：}\videoduration\par
\textbf{视频链接：}
\ifdefempty{\videourl}{未填写}{\href{\videourl}{\nolinkurl{\videourl}}}\par
\textbf{输入模式：}\inputmode\par
\textbf{处理说明：}\processingnote
\end{tcolorbox}
\end{titlepage}

\tableofcontents
\newpage

%% Replace this block with generated Chinese teaching notes.
%%
%% Required body structure:
%% - for trading videos, default to a zero-foundation textbook, not a compressed summary
%% - include or reference source_coverage_report before finalizing
%% - follow writing_plan.json instead of dumping subtitle lines
%% - each major section should explain motivation, problem/context, core idea,
%%   mechanism, evidence/example, support boundary, and learner takeaway when applicable
%% - substantial chapters should include learning goals, prerequisites, source scene,
%%   timestamped source evidence, novice explanation, mechanism breakdown,
%%   common misconceptions, risk boundary, chapter summary, and review answers
%% - define specialist terms before relying on them
%% - define each first-use trading term with plain language, source scene,
%%   professional meaning, common misunderstanding, and misuse boundary
%% - explain intuition before formulas, code, or trading jargon
%% - use full trade-case replay for concrete trades, cancelled plans, stop moves,
%%   partial exits, exits, and no-trade decisions
%% - include selected source frames with \sourceframefig
%% - include crops with \sourceframefig
%% - include AI-generated teaching diagrams with \teachingdiagram
%% - introduce every figure before insertion and interpret its evidence boundary after insertion when needed
%% - keep images outside boxes
%% - use importantbox/knowledgebox/warningbox only for high-signal teaching payloads
%% - end each major section with \subsection{本章小结}
%% - end the document with a final synthesis section covering core claims,
%%   cross-section links, practical implications, limits, and open questions when supported
%%
%% Example:
%% \section{示例章节}
%% 正文说明。
%% \begin{importantbox}{核心概念}
%% 重要解释。
%% \end{importantbox}
%% \sourceframefig{figures/fig001.jpg}{截图说明}{00:01:20--00:01:28}{0.9}
%% \teachingdiagram{diagrams/diagram001.svg}{教学图说明}{KP001，00:01:20--00:01:28}{0.78}
%% \subsection{本章小结}
%% 小结。

\end{document}
